% ------------------------------------------------------------------ at a glance
\clearpage
\thispagestyle{plain}
\section*{At a Glance}
\addcontentsline{toc}{section}{At a Glance}

{\large\itshape\color{teal} Remoteness is no longer a fixed cost paid once.
It is a recurring tax, and the taxes compound.\par}
\vspace{14pt}

\noindent
\begin{minipage}[t]{0.32\linewidth}
\begin{claimtile}{1\enspace Far, not small}
Distance carries the income penalty. Conditional on governance, doubling
economic remoteness goes with \regDoubleRemotePctAbs\% lower income per head;
population barely matters (\cref{sec:edge}).
\end{claimtile}
\end{minipage}\hfill
\begin{minipage}[t]{0.32\linewidth}
\begin{claimtile}{2\enspace Arrival converges, intensity doesn't}
New technologies reach the edge within about a decade, not \lagOldMedian{} years.
But broadband, bandwidth and redundancy still lag: \nSIDSOneCable{} island
states hang on one cable or none (\cref{sec:tech,sec:accel}).
\end{claimtile}
\end{minipage}\hfill
\begin{minipage}[t]{0.32\linewidth}
\begin{claimtile}{3\enspace The edge is not destiny}
Iceland, the Faroes, Singapore and Mauritius escaped through rents,
institutions and becoming hubs. But even escapees export talent: Mauritius
lost \emigMauritius\% of its graduates (\cref{sec:escape}).
\end{claimtile}
\end{minipage}

\vspace{18pt}
\noindent{\sffamily\bfseries\small\color{teal}Headline chart\enspace}{\sffamily\small\color{slate}Access versus intensity, medians per 100 people}
\par\smallskip
\begin{center}
\begin{tikzpicture}
\begin{axis}[edge, width=0.92\linewidth, height=0.42\linewidth,
    xmin=1998, xmax=2023, xtick={2000,2005,2010,2015,2020},
    x tick label style={/pgf/number format/1000 sep={}},
    ymin=0, ymax=100,
    legend style={at={(0.02,0.98)},anchor=north west}, legend columns=2]
  \addplot[lineSlate] table[x=year,y=net_hic] {\datadir/internet_groups.csv};
  \addlegendentry{Online, rich}
  \addplot[lineTeal] table[x=year,y=net_sids] {\datadir/internet_groups.csv};
  \addlegendentry{Online, islands}
  \addplot[lineSlate, dashed, line width=1.1pt] table[x=year,y=bb_hic] {\datadir/internet_groups.csv};
  \addlegendentry{Broadband, rich}
  \addplot[lineEmber] table[x=year,y=bb_sids] {\datadir/internet_groups.csv};
  \addlegendentry{Broadband, islands}
\end{axis}
\end{tikzpicture}
\end{center}
{\sffamily\footnotesize\color{slate}
Small island states have nearly caught up on being online
(\netSidsLast\% vs \netHicLast\% in \windowEnd) but not on fixed broadband
(\bbSidsLast{} vs \bbHicLast{} lines per 100). Source: World Bank WDI; see
\cref{fig:internet}.\par}

\vspace{14pt}
\begin{counterpoint}[What the data do not show]
Remote places have not, in general, become poorer: most gained in absolute
terms. Several gaps have \emph{narrowed} (internet access; shipping relative
to the old hubs). The claims here are about relative position, volatility and
control, and the paper flags where the evidence cuts against them.
\end{counterpoint}
\clearpage
